\begin{table}[H]
\centering\footnotesize
\caption{Historical O2C Portfolio Performance}\label{tab:horse_race}
\begin{threeparttable}
\renewcommand{\arraystretch}{1.12}
\begin{tabular*}{\textwidth}{@{\extracolsep{\fill}}lrrrrr@{}}
\toprule
Signal & Mean (bps) & Annual (\%) & Sharpe & CH4 $\alpha$ & U.S. q5 $\alpha$ \\
\midrule
ESI (EW) & 5.31 (5.95) & 12.85 & 1.54 & 5.17 (5.74) & 5.42 (6.06) \\
Investor tone (EW) & 4.43 (4.22) & 10.73 & 1.11 & 4.73 (4.55) & 4.60 (4.34) \\
Manager tone (EW) & 2.48 (2.83) & 6.00 & 0.71 & 2.36 (2.65) & 2.46 (2.80) \\
ESI (VW) & 5.26 (3.57) & 12.72 & 0.91 & 5.16 (3.46) & 5.40 (3.67) \\
Investor tone (VW) & 4.38 (2.70) & 10.59 & 0.69 & 4.73 (2.88) & 4.51 (2.79) \\
Manager tone (VW) & 2.75 (1.79) & 6.66 & 0.45 & 2.73 (1.77) & 2.69 (1.75) \\
\bottomrule
\end{tabular*}
\begin{tablenoteblock}
Gross long-minus-short returns; 3,642 dates for each row. Parentheses contain HAC(5) $t$ statistics. Up to 50 positive-signal stocks in the upper 30 percent form the long leg; up to 50 negative-signal stocks in the lower 30 percent form the short leg. Both legs must exist. Annualization is the historical arithmetic convention, $242\bar r\times100$, with Sharpe scaled by $\sqrt{242}$. Candidate pools are signal-specific, not a common-stock-pool horse race. Next-record O2C excess returns, future closing-limit exclusions, full-period text scales, and signal-date factors reproduce the original implementation. CH4 is domestic; q5 is the U.S. Global-q series. The intraday outcome and daily factor intervals differ. These are historical diagnostic alphas, not implementable profits. All figures are recomputed values rather than copied PDF numbers.
\end{tablenoteblock}
\end{threeparttable}
\end{table}
